% Pursuing Stacks (cotes 134-2 to 134-8): archive cote 134-4, site page 18 (PDF page 19), typescript page 106.
% typescript. Draft compared: ps3.tex:886-940.
% AI-transcribed by Claude Opus 5.5 (claude-opus-5-5, adaptive thinking, effort high), Message Batches API, batch ps1 (msgbatch_01Rv5hCaNTtUt4JHx1yqoHPL), 2026-11-12; image review recorded in records/CHECK_DECISIONS_134_3_4.tsv.
\scanpage{134-4/18}
\mspage{106}
set of objects of a category ...).

We now get canonical functors
\[ M\hand{\longrightarrow} M_{\underline{I}} \hand{\to} M_W \hand{\to} H \quad , \]
where H can again be considered as the common localization of the three
categories M, $M_{\underline{I}}$ (any $\underline{I}$), $M_W$ with respect to the notion of weak equivalen%
ce in each, It is clear that the set of weak equivalences in $M_W$ , say $\overset{\hand{\rule{0.65em}{0.4pt}}}{W}$,
is saturated provided W is. Also, the canonical functor $M\hand{\to} M_{\over[typed]{\illegible{1}}{W}}$ commutes
with finite products, and also with finite sums provided formation of
such sums in M commutes with taking products \over{w}{w}\over{\unsure{i|j}}{i}\over{t}{t}h \deleted{a}\inserted[typed]{any} fixed element of M.
\deleted{T\over{h}{h}e same statement will hold for infinite products or sums, provided we
assume there exists a characteristic homotopy interval \inserted[typed]{$\underline{I}_o$}, in M, so that we
got a description
\[ M_W \qquad M_{\underline{I}_o} \quad , \]
and may apply the \xout{result} corresponding fact for M \quad $M_{\underline{I}_o}$ .}

A map f: $X\hand{\to} Y$ is called a \inserted{\uline{homotopy equivalence} or} \uline{homotopism} (with respect to W) if it is
an isomorphism in $M_W$ , namely if there exists g: $Y\hand{\to} X$ such that
\[ gf \underset{W}{\hand{\sim}} \mathrm{id}_X \quad , \quad fg \underset{W}{\hand{\sim}} \mathrm{id}_Y \quad . \]
\inserted[typed]{This implies f and g are weak equivalences.}
\xout{When X\illegible{1}simpl}\hand{\ensuremath{\unsure{\lozenge}}} If such an f exists namely if X and Y are isomorphic
objects of $M_{\over{W}{W}}$ , we'll say they are \inserted[typed]{W-}\uline{homotopic}, or simply \uline{homotopic}. This
implies that there \xout{are}\inserted[typed]{exist} weak equivalences $X\hand{\to} Y$ and $Y\hand{\to} X$, but the con-%
verse of course is'nt always true. If Y is the final object $e_M$ , we'll
say X is W-\uline{contractible} (or simply \uline{contactible}) \xout{in order} instead of
W-homotopic to Y\hand{$=e$.} In case there exists a characteristic homotopy
interval $\underline{I}_o$ , all these W-notions boil down to the corresponding $\underline{I}_o$-noti-%
ons considered before.

\inserted{\uline{Remark.}} It occurs to me that in all what precedes, we never made any use
really of the existence of a homotopy interval\over{\unsure{.}}{!} T\over[typed]{H}{h}e only notion we have
been effectively been working with, it seems to me, is the notion\unsure{s} of
``weak homotopy interval'', by which I mean a\xout{n object} triple $(\underline{I},\hand{\ensuremath{\delta}}_o,\hand{\ensuremath{\delta}}_1)$
satisfying merely the condition \xout{a) and (\illegible{1}) of page 59} that I be aspheric
over e (which is the condition c) on page 59). Such an \underline{I} always exists
of course, we need only take I = e itself ! I\over[typed]{f}{n} \inserted[typed]{case} however an \inserted{actual} homotopy
interval (with Ker$(\hand{\ensuremath{\delta}}_o,\hand{\ensuremath{\delta}}_1) = \hand{\varnothing}_M$ ) does exist, to make sure that the notion
of W-homotopy obtained by using all weak homotopy intervals is the same
as when using only actual homotopy intervals, we should be sure that
for any I aspheric over e, any two sections $\hand{\ensuremath{\delta}}_o,\hand{\ensuremath{\delta}}_1$ of I over e are W-homo-%
topic in the initial meaning. This is evidently so when M = (Cat) and
W = $\hand{\ensuremath{\mathcal{W}}}_{(Cat)}$ (indeed, it is enough that I be O-connected, instead of

\margin{18}
